\documentclass[11pt]{article}
\usepackage[margin=2.5cm]{geometry}
\usepackage{amsmath}
\usepackage{booktabs}
\usepackage{parskip}

\title{The cloglog link, simply}
\date{}

\begin{document}
\maketitle

\section*{The problem a link solves}

The model has to give each segment a probability $p$ of developing plaque
between the baseline and the follow-up scan. A probability must lie between
0 and 1, but a regression $\eta = \beta_0 + \beta_1 x_1 + \dots$ can take any
value. A link function turns $p$ into a number on the whole real line, so
that the regression can be written on that scale and turned back into a
probability afterwards.

Logistic regression uses the logit link, $\log\{p/(1-p)\}$, the log odds.
Its coefficients are odds ratios. The complementary log-log (cloglog) link is
\[
  \operatorname{cloglog}(p) = \log\{-\log(1-p)\} = \eta .
\]

\section*{Where cloglog comes from}

Think of plaque as something that can start at any moment during the ten
years, at some rate, the hazard. Adding the hazard up over the follow-up
gives the cumulative hazard $H$. If the risk builds up this way, the chance
of still being free of plaque at the follow-up scan is
\[
  1 - p = e^{-H}.
\]
Taking logs twice undoes this: $\log\{-\log(1-p)\} = \log H$. The cloglog
link is therefore nothing more than the log of the cumulative hazard, and the
model says
\[
  \log H = \beta_0 + \beta_1 x_1 + \dots
\]
Raising a feature by one unit multiplies $H$ by $e^{\beta_1}$. This is a
hazard ratio, the same quantity a Cox model reports, even though we never see
when the plaque started, only that it was absent at baseline and present at
follow-up.

\section*{Why it suits the segment model}

A person is free of plaque only if all nine segments are. Under cloglog
each segment has $1-p_s = e^{-H_s}$, so
\[
  P(\text{no plaque in any segment}) = e^{-H_1}\cdots e^{-H_9}
  = e^{-(H_1+\dots+H_9)} .
\]
The segment hazards simply add up to one person hazard, and the person risk
is one minus this. With the person random intercept $b$, every segment hazard
is multiplied by the same $e^{b}$, and the person risk is found by averaging
over $b$. The logit link has no such simple sum, which is the main reason
cloglog was chosen.

\section*{How different is it from logit?}

When the risk is small, odds ratios and hazard ratios are almost equal, and
the two links give nearly the same answer. They part when the risk is large.
For a hazard ratio of 1.5:

\begin{center}
\begin{tabular}{lcc}
\toprule
Risk in the reference group & Odds ratio & Level \\
\midrule
2\%  & 1.51 & segment \\
15\% & 1.56 & segment \\
42\% & 1.75 & person \\
\bottomrule
\end{tabular}
\end{center}

At segment level the choice therefore matters little for the coefficients,
the likelihood-ratio test of the geometry block or the C-index. It matters
mainly for the pooled person risk and its calibration.

\section*{What to keep in mind}

The hazard-ratio reading rests on proportional hazards, which cannot be
tested with only one follow-up scan. The hazard ratios are conditional on the
person random intercept, so they describe a change within the same person's
risk level, not an average over the population. In the current code
(\texttt{CGPS\_data/progression\_who\_pooled.py}) log follow-up time enters
as a covariate with its own coefficient, not as a fixed offset; a coefficient
of 1 would mean a constant hazard over time.

\end{document}
